% !PRINTED 164
\weberpage{164}
If the group $S$ of degree $\mu$ has an invariant of index
 $\mu$, then $T$ becomes the unit group and $S$ is a cyclic
group, which consists of the elements $1, E, E^2 \ldots E^{\mu-1}$.

If we consider as an example the symmetric permutation group
 $P$ of $n$ elements $x_1,x_2,\ldots,x_n$, which by \S.\,39 is indeed contained
among the groups of linear substitutions, then we have
as absolute invariants the \gesperrt{symmetric functions} of the
$n$ variables $x$. The difference product
\[
\sqrt{\varDelta} = (x_1-x_2)(x_1-x_3)\ldots(x_{n-1}-x_n)
\]
is a relative invariant of index 2, to which the alternating
group belongs. Since the group $P$ has no normal divisor
apart from the alternating group and is also not cyclic,
there are no relative invariants of index higher than 2.

The absolute invariants, i.e. the symmetric functions,
are here rationally representable by a finite number
of such forms, namely by the symmetric fundamental functions,
and the invariants of index 2 are the product
of $\sqrt{\varDelta}$ with absolute invariants. We shall prove below that analogous theorems also
hold in the general case.

We close here with the proof of a general
theorem which in all applications is of great use for the formation
of the invariants of a group $S$.

In \S.\,60 of the first volume we defined, for an arbitrary form
of the $n$ variables $F(x_1,x_2,\ldots x_n) = F(x)$, certain forms of the same
 variables $C(x_1,x_2,\ldots x_n) = C(x)$ as covariants,
which were characterised by the property that, when $F(x)$ is transformed
by an arbitrary linear substitution into a new form $F'(y)$,
 the form $C$ satisfies the condition
\[
C'(y) = r^\lambda\,C(x),
\]
where $C'$ depends on the coefficients of $F'$ just as $C$ does on the
coefficients of $F$. Here $r$ denotes the substitution determinant,
and is thus a constant. The coefficients of the
form $F$ occur in $C$ only homogeneously. As an example of a covariant
 we cite the \gesperrt{Hessian determinant}.

If now $F(x)$ is an invariant of the group $S$, and the
$y$ are likewise connected with the $x$ by a substitution from $S$,
then the coefficients of $F$ differ only by
a common constant factor from the coefficients
of $F'$, and the same therefore also holds of the two forms $C$
%CONT>>
